\section*{Chapter \ref{ch:b2:enolates-aldol} --- Enols, Enolates and the Aldol Reaction}

\begin{solution}{exo:b2:enolates-aldol:1}
Butanone: but-1-en-2-ol \ce{CH2=C(OH)CH2CH3} and but-2-en-2-ol \ce{CH3C(OH)=CHCH3} ((\textit{E}) and
(\textit{Z})). Cyclohexanone: cyclohex-1-en-1-ol.
\end{solution}

\begin{solution}{exo:b2:enolates-aldol:2}
Ethane $<$ propanone $<$ ethyl 3-oxobutanoate (10.7) $<$ pentane-2,4-dione (9.0).
\end{solution}

\begin{solution}{exo:b2:enolates-aldol:3}
3-Hydroxybutanal, then but-2-enal on heating.
\end{solution}

\begin{solution}{exo:b2:enolates-aldol:4}
Diethyl propanedioate, sodium ethoxide, 1-bromobutane; hydrolysis to butylpropanedioic acid;
heating removes \ce{CO2}: \ce{CH3(CH2)3CH2COOH}, hexanoic acid.
\end{solution}

\begin{solution}{exo:b2:enolates-aldol:5}
Kinetic: deprotonation at C6 (the \ce{CH2} side), with LDA at \qty{-78}{\degreeCelsius}. Thermodynamic:
the more substituted double bond C1=C2, towards the methyl-bearing carbon, with an alkoxide in its
alcohol at room temperature.
\end{solution}

\begin{solution}{exo:b2:enolates-aldol:6}
4-Phenylbut-3-en-2-one (benzalacetone), \ce{C6H5CH=CHCOCH3}. Benzaldehyde has no $\alpha$ hydrogen: it
cannot form an enolate and is only the electrophile.
\end{solution}

\begin{solution}{exo:b2:enolates-aldol:7}
The enolate of one methyl group (C1) attacks the other carbonyl (C5): a five-membered ring,
3-methylcyclopent-2-en-1-one after dehydration. The alternative, the enolate at C3 attacking C5,
would make a strained three-membered ring.
\end{solution}

\begin{solution}{exo:b2:enolates-aldol:8}
\ce{2CH3COOC2H5 -> CH3COCH2COOC2H5 + C2H5OH} with sodium ethoxide; the product, deprotonated, is
recovered as ethyl 3-oxobutanoate on acidification.
\end{solution}

\begin{solution}{exo:b2:enolates-aldol:9}
Propanone, ethanal and acetophenone (they have a \ce{CH3-CO} group; ethanal is \ce{CH3CHO}).
Pentan-3-one does not.
\end{solution}

\begin{solution}{exo:b2:enolates-aldol:10}
The enolate at one end attacks the ester at the other: ethyl 2-oxocyclopentane-1-carboxylate, a
five-membered ring.
\end{solution}

\begin{solution}{exo:b2:enolates-aldol:11}
In 3-methylpentan-2-one the stereocentre (C3) is the $\alpha$ carbon: enolisation makes it planar, and
reprotonation from either face racemises it. In 4-methylhexan-2-one the stereocentre is at C4,
$\beta$ to the carbonyl, untouched by enolisation.
\end{solution}

\begin{solution}{exo:b2:enolates-aldol:12}
Propanone with two equivalents of benzaldehyde in base. With two different aldehydes, condense
propanone with one aldehyde first (one equivalent, isolating benzalacetone), then condense its
remaining methyl group with the second. In one pot the two aldehydes compete and give a mixture of
three dienones.
\end{solution}

\begin{solution}{pb:b2:enolates-aldol:1}
\textbf{1.} Propanone, six $\alpha$ hydrogens (three on each methyl).
\textbf{2.} $\mathrm{CH_3{-}CO{-}CH_2^-} \leftrightarrow \mathrm{CH_3{-}C(O^-){=}CH_2}$.
\textbf{3.} It has no hydrogen on the carbon next to its carbonyl (that carbon belongs to the ring and
carries no H).
\textbf{4.} Benzaldehyde, an aldehyde in excess, is the better electrophile; propanone's enolate meets it
more often than a propanone carbonyl.
\textbf{5.} A ketone has two carbon groups pushing electron density onto its carbonyl carbon and
hindering approach; an aldehyde has one.
\textbf{6.} Only a small fraction of enolate is needed: it is regenerated as it reacts.
\textbf{7.} The enolate carbon bonds to the carbonyl carbon of benzaldehyde: the alkoxide of
4-hydroxy-4-phenylbutan-2-one.
\textbf{8.} Removal of an $\alpha$ hydrogen and loss of hydroxide from the $\beta$ carbon: 4-phenylbut-3-en-2-one.
\textbf{9.} The double bond formed is conjugated with both the ring and the carbonyl.
\textbf{10.} The other methyl group still has three $\alpha$ hydrogens.
\textbf{11.} \ce{2C6H5CHO + CH3COCH3 -> C6H5CH=CHCOCH=CHC6H5 + 2H2O}.
\textbf{12.} The dehydrations, and the precipitation of the product, which leaves the solution.
\textbf{13.} Two, both (\textit{E}).
\textbf{14.} The (\textit{Z}) isomers put the phenyl group next to the carbonyl group: steric strain.
\textbf{15.} Two \ce{C=C}, one \ce{C=O} and the two rings: a conjugated system across the molecule.
\textbf{16.} Its extended conjugation narrows the $\pi$ gap (\cref{prop:b2:huckel:gap}): it absorbs at the
blue end of the visible and looks yellow.
\textbf{17.} No: it is planar (or nearly) with a mirror plane.
\textbf{18.} 106.12, 58.08 and \qty{234.30}{g/mol}.
\textbf{19.} Benzaldehyde $2.1 \times 1.045/106.12 = \qty{20.7}{mmol}$; propanone $0.75 \times
0.7845/58.08 = \qty{10.1}{mmol}$.
\textbf{20.} Propanone needs \qty{20.3}{mmol} of benzaldehyde, and \qty{20.7}{mmol} is available:
propanone is limiting, barely.
\textbf{21.} $0.01013 \times 234.30 = \qty{2.37}{g}$.
\textbf{22.} Leftover benzaldehyde stays in the ethanol and is washed away.
\textbf{23.} The mono-condensation product, benzalacetone.
\textbf{24.} $0.75 \times 2.37 = \boldsymbol{\approx \qty{1.78}{g}}$ of dibenzalacetone.
\end{solution}
